\section{Discussion}
\label{sec:discussion}
\looseness=-1 We introduce a constrained learning framework for first-order debiasing in causal estimation and semiparametric inference.
We pose asymptotically optimal plug-in estimators as those whose nuisance parameters are solutions to a optimization problem, under
the constraint that the first-order error of the plug-in estimator with respect to the nuisance parameter estimate is zero. 
This perspective encompasses versions of one-step estimation and targeting. %

\looseness=-1 The constrained learning perspective enables %
a new method %
(Constrained Learner, a.k.a. C-Learner), which solves this constrained optimization directly while using the entire model class. 
It outperforms existing asymptotically optimal methods without additional heuristics or assumptions in settings with low overlap, and performs similarly otherwise. 
We demonstrate C-Learner's versatility by instantiating it with model classes including linear models, gradient boosted trees, and neural networks, and on datasets with both tabular and text covariates. 
We also construct a theoretical example to understand how C-Learner may achieve stable estimates in settings with low overlap: we assume inverse propensity weights are heavy tailed and outcome models lie within a square-integrable envelope. Then, C-Learner has finite variance while basic versions of one-step estimation and targeting do not. 

\looseness=-1 Our theoretical analysis is only a small initial step in building a principled understanding of the benefits of the constrained learning framework. One future direction is to investigate which model classes and optimization methods satisfy our theoretical assumptions (e.g. Assumption \ref{ass:empirical_proc_mean}).
Finally, 
we hope this work spurs further investigation on how constrained optimization can enable more robust estimators, by proposing better optimization procedures, or extending to additional model classes and estimands beyond \Cref{sec:other_estimands}. 


